\section{Method}
\label{sec:method}

% Describe what the code in sumo_rl/agents/clm_agent.py actually does; re-check the code before editing.

\subsection{Problem Formulation}
% Signal control as a sequential decision problem: at each decision step (every delta_time seconds, subject to
% min/max green and yellow), the controller picks the next green phase.
\todo[inline]{Formulation: decision interval, action set, constraints.}

\subsection{Describing the Intersection as Text}
% describe_state / lane_labels: how lanes, queues, waiting times and the current phase become a text state.
% Include one real example state string (from a run with -v) in a figure or listing.
\todo[inline]{State description, with an example.}

\subsection{Describing Phases as Candidate Actions}
% describe_phases: each green phase as a text description of its movements (approach + turn).
\todo[inline]{Phase descriptions, with an example.}

\subsection{Ranking Phases with a CLM}
% Query (DEFAULT_QUESTION) + state vs. candidate phases; score = dot product of embeddings; pick the argmax.
% Caching of phase embeddings and the client-side LRU cache.
\todo[inline]{Ranking and caching.}

\subsection{Traffic Specialization}
% clm_collect_labels.py: lookahead labelling in a shadow simulation, positive = best phase over the horizon.
% Fine-tuning of the projection heads only (backbone frozen). Labels from the train split only.
\todo[inline]{Fine-tuning procedure, once the teacher/label source is decided.}
